% part: intuitionistic-logic
% chapter: soundness-completeness
% section: lindenbaum

\documentclass[../../../include/open-logic-section]{subfiles}

\begin{document}

\olfileid{int}{sc}{lin}

\olsection{લિન્ડનબાઉમનું સહાયક પ્રમેય}

સ્ફુરણાવાદી તર્ક માટેનો પૂર્ણતા પ્રમેય $\Gamma \Proves/ !A$
માનીને $\mSat{M}{\Gamma}$ અને~$\mSat/{M}{!A}$ ધરાવતું
નિદર્શ રચીને સાબિત થાય છે.

પ્રશિષ્ટ તર્કમાં !!{derivability}ના સંબંધને સુસંગતતાના
ખ્યાલમાં ઉતારી શકાય છે: !!a{formula}~$!A$ એ
!!{formula}sના ગણમાંથી ત્યારે અને માત્ર ત્યારે જ
!!{derivable} હોય છે જ્યારે તે ગણ અને~$!A$નો નિષેધ
મળીને અસુસંગત થાય. સ્ફુરણાવાદી તર્કમાં આવું શક્ય નથી.
ત્યાં જો $\lnot!A$ અસુસંગત હોય, તો માત્ર
$\Proves \lnot\lnot !A$ મળે છે. સામાન્ય રીતે
$\lnot\lnot!A \lif !A$ સ્ફુરણાવાદી રીતે માન્ય નથી,
તેથી $\Proves !A$ તારવી શકાતું નથી.

આથી નિદર્શ~$\mModel{M}$ રચતી વખતે !!{formula}~$!A$નું
!!{derivability} ન હોવું ધ્યાનમાં રાખવું પડે છે. તેથી નિદર્શ
$\mModel{M}$ રચવા માટે પૂર્ણ ગણ $\Gamma^* \supseteq \Gamma$
વાપરી શકાતો નથી, કારણ કે દરેક પૂર્ણ ગણ~$\Gamma^*$માં
$\Gamma^* \Proves !A \lor \lnot !A$ હોય છે.

પૂર્ણ ગણ~$\Gamma^*$ને બદલે આપણે સૂત્રોના અભાજ્ય ગણનો
ખ્યાલ વાપરીશું:

\begin{defn}\ollabel{defn:prime}
  !!{formula}sનો ગણ~$\Gamma$ ત્યારે અને માત્ર ત્યારે જ
  \emph{અભાજ્ય} કહેવાય જ્યારે
  \begin{enumerate}
  \item\ollabel{defn:prime1} $\Gamma$ સુસંગત હોય, એટલે
    $\Gamma \Proves/ \lfalse$;
  \item\ollabel{defn:prime2} જો $\Gamma \Proves !A$ હોય,
    તો $!A \in \Gamma$; અને
  \item\ollabel{defn:prime3} જો $!A \lor !B \in \Gamma$
    હોય, તો $!A \in \Gamma$ અથવા $!B \in \Gamma$.
  \end{enumerate}
\end{defn}

\begin{lem}[લિન્ડનબાઉમનું સહાયક પ્રમેય]
  \ollabel{lem:lindenbaum} જો $\Gamma \Proves/ !A$ હોય,
  તો એવો $\Gamma^* \supseteq \Gamma$ છે કે $\Gamma^*$
  અભાજ્ય છે અને $\Gamma^* \Proves/ !A$.
\end{lem}

\begin{proof}
  $!B_1 \lor !C_1$, $!B_2 \lor !C_2$, \dots, એ
  $!B \lor !C$ સ્વરૂપના બધા !!{formula}sની પરિગણના હોય.
  !!{formula}sના ગણોની વધતી શ્રેણી~$\Gamma_n$ વ્યાખ્યાયિત
  કરીએ છીએ, જેમાં દરેક $\Gamma_{n+1}$ એ $\Gamma_n$માં
  એક નવું !!{formula} ઉમેરીને બને છે. બધા~$\Gamma_n$નો
  યોગગણ $\Gamma^*$ હશે. નવા !!{formula}s એવી રીતે પસંદ
  કરીએ છીએ કે $\Gamma^*$ અભાજ્ય રહે અને છતાં
  $\Gamma^* \Proves/ !A$ હોય. એટલે દરેક તબક્કે
  નીચેની શરતો સંતોષતું પહેલું વિયોજન~$!B_i \lor !C_i$
  શોધવું પડશે:
  \begin{enumerate}
  \item\ollabel{gamma-1} $\Gamma_n \Proves !B_i \lor !C_i$
  \item\ollabel{gamma-2} $!B_i \notin \Gamma_n$ અને
    $!C_i \notin \Gamma_n$
  \end{enumerate}
  જો $\Gamma_n \cup \{!B_i\} \Proves/ !A$ હોય, તો
  $\Gamma_n$માં $!B_i$ ઉમેરીએ; નહિ તો $!C_i$ ઉમેરીએ.
  આ રીત કામ કરે છે તે બતાવવું પડશે. હાલ માટે
  \olref{gamma-1} અને~\olref{gamma-2} સંતોષતા સૌથી
  નાના $i$ને $i(n)$ કહીએ.

  $\Gamma_0 = \Gamma$ મૂકો અને
  \[
  \Gamma_{n+1} = \begin{cases}
    \Gamma_n \cup \{!B_{i(n)}\} &
    \text{જો $\Gamma_n \cup \{!B_{i(n)}\} \Proves/ !A$ હોય} \\
    \Gamma_n \cup \{!C_{i(n)}\} & \text{અન્યથા}
    \end{cases}
  \]
  જો $i(n)$ વ્યાખ્યાયિત ન હોય, એટલે $\Gamma_n \Proves !B \lor !C$
  થાય ત્યારે હંમેશાં $!B \in \Gamma_n$ અથવા $!C \in \Gamma_n$
  હોય, તો $\Gamma_{n+1} = \Gamma_n$ મૂકો.
  હવે $\Gamma^* = \bigcup_{n=0}^\infty \Gamma_n$ મૂકો.

  પહેલાં દરેક $n$ માટે $\Gamma_n \Proves/ !A$ બતાવીએ.
  $n$ પર આગમન કરીએ. $n = 0$ માટે સહાયક પ્રમેયની ધારણા,
  એટલે $\Gamma \Proves/ !A$, પરથી દાવો મળે છે.
  આગમનના પગલામાં કોઈ પણ $n\ge 0$ માટે
  $\Gamma_n \Proves/ !A$ પરથી
  $\Gamma_{n+1} \Proves/ !A$ બતાવવું છે. $i(n)$
  વ્યાખ્યાયિત ન હોય, તો $\Gamma_{n+1} = \Gamma_n$
  હોવાથી કશું સાબિત કરવાનું નથી. તેથી $i(n)$
  વ્યાખ્યાયિત છે એમ ધારો. સરળતા માટે $i = i(n)$ મૂકો.
  \sourcecorrection{OLMD-093}{મૂળમાં આધાર $n=0$ પછી
  આગમનનું પગલું $n>0$થી શરૂ થાય છે, જેથી
  $\Gamma_0$થી $\Gamma_1$નો કિસ્સો છૂટી જાય છે.
  પગલું બધા $n\ge 0$ માટે જણાવ્યું છે.}

  દાવાનું પ્રતિપક્ષી વિધાન સાબિત કરીએ. ધારો કે
  $\Gamma_{n+1} \Proves !A$. રચના મુજબ જો
  $\Gamma_n \cup \{!B_i\} \Proves/ !A$ હોય, તો
  $\Gamma_{n+1} = \Gamma_n \cup \{!B_i\}$; નહિ તો
  $\Gamma_{n+1} = \Gamma_n \cup \{!C_i\}$.
  પહેલો કિસ્સો બની શકે નહીં, કારણ કે ત્યારે
  $\Gamma_{n+1} \Proves/ !A$ થાય. તેથી
  $\Gamma_n \cup \{!B_i\} \Proves !A$ અને
  $\Gamma_{n+1} = \Gamma_n \cup \{!C_i\}$.
  $i(n)$ની વ્યાખ્યાથી $\Gamma _n \Proves !B_i \lor !C_i$.
  વળી $\Gamma_n \cup \{!B_i\} \Proves !A$ અને
  $\Gamma_{n+1} = \Gamma_n \cup \{!C_i\} \Proves !A$.
  આથી ઇચ્છિત $\Gamma_n \Proves !A$ મળે છે.

  જો $\Gamma^* \Proves !A$ હોય, તો એવો સાન્ત ઉપગણ
  $\Gamma' \subseteq \Gamma^*$ હોય
  કે $\Gamma' \Proves !A$. તેનું દરેક
  $!D \in \Gamma'$ કોઈક~$i$ માટે $\Gamma_i$માં હશે.
  આ સૂચકોમાં સૌથી મોટો $n$ લો. $i \le n$ હોય ત્યારે
  $\Gamma_i \subseteq \Gamma_n$ હોવાથી
  $\Gamma' \subseteq \Gamma_n$. તેથી
  $\Gamma_n \Proves !A$, જે ઉપર સાબિત કરેલા
  $\Gamma_n \Proves/ !A$થી વિરુદ્ધ છે.

  છેલ્લે $\Gamma^*$ અભાજ્ય છે, એટલે કે
  \olref{defn:prime}ની \olref{defn:prime1},
  \olref{defn:prime2} અને~\olref{defn:prime3}
  શરતો સંતોષે છે તે બતાવીએ.

  સૌપ્રથમ $\Gamma^* \Proves/ !A$ છે, તેથી
  $\Gamma^*$ સુસંગત છે અને \olref{defn:prime1} મળે છે.

  હવે જો $\Gamma^* \Proves !B \lor !C$ હોય, તો
  $!B \in \Gamma^*$ અથવા $!C \in \Gamma^*$ છે તે
  બતાવીએ. આનાથી \olref{defn:prime3} મળે છે, કારણ કે
  $!B \lor !C \in \Gamma^*$ હોય, તો
  $\Gamma^* \Proves !B \lor !C$ પણ હોય. તેથી
  $\Gamma^* \Proves !B \lor !C$, પરંતુ
  $!B \notin \Gamma^*$ અને $!C \notin \Gamma^*$,
  એમ ધારો. $\Gamma^* \Proves !B \lor !C$
  હોવાથી કોઈક~$n$ માટે $\Gamma_n \Proves !B \lor !C$.
  બધા વિયોજનોની પરિગણનામાં $!B \lor !C$ કોઈક
  $!B_j \lor !C_j$ તરીકે આવે છે. $!B_j \lor !C_j$
  એ $i(n)$ની વ્યાખ્યાની શરતો સંતોષે છે: એટલે
  $\Gamma_n \Proves !B_j \lor !C_j$, જ્યારે
  $!B_j \notin \Gamma_n$ અને $!C_j \notin \Gamma_n$.
  દરેક તબક્કે પસંદ કરેલું વિયોજન~$!B_i \lor !C_i$
  પછીની યાદીમાં ફરી શરતો સંતોષતું નથી, કારણ કે
  તેમાં $!B_i$ અથવા $!C_i$ ઉમેરાય છે. આ સૂચક
  પહેલાં સાન્ત સૂચકો જ છે; તેથી આ વિયોજન શરતો સંતોષતું
  રહે ત્યાં સુધી તેનાથી નાના બાકી સૂચકો એક પછી એક
  દૂર થશે અને કોઈક તબક્કે~$m$ માટે $j = i(m)$ મળશે.
  \sourcecorrection{OLMD-094}{મૂળ કહે છે કે દરેક તબક્કે
  શરતો સંતોષતા વિયોજનોની કુલ સંખ્યા ઓછામાં ઓછી એકથી
  ઘટે છે; નવા વિયોજનો નિગમ્ય બની શકે હોવાથી આ દાવો
  જરૂરી નથી. નક્કી કરેલા $j$થી નાના સાન્ત સૂચકો
  ફરી પસંદ ન થાય એ તર્ક વાપર્યો છે.}
  પરંતુ પછી $!B \in \Gamma_{m+1}$ અથવા
  $!C \in \Gamma_{m+1}$, જે $!B \notin \Gamma^*$
  અને $!C \notin \Gamma^*$ની ધારણાથી વિરુદ્ધ છે.

  હવે ધારો કે $\Gamma^* \Proves !B$. તો
  $\Gamma^* \Proves !B \lor !B$. હમણાં સાબિત કર્યું
  છે કે $\Gamma^* \Proves !B \lor !B$ હોય, તો
  $!B \in \Gamma^*$. તેથી $\Gamma^*$ એ
  \olref{defn:prime}ની \olref{defn:prime2} સંતોષે છે.
\end{proof}

\begin{prob}
જો $\Gamma \Proves/ \bot$ હોય, તો પ્રશિષ્ટ તર્કમાં
$\Gamma$ સુસંગત છે, એટલે~$\Gamma$નાં બધાં સૂત્રોને
સાચાં બનાવતું મૂલ્યાંકન છે, તે બતાવો.
\end{prob}

\end{document}
